\documentclass[preview]{standalone}

\usepackage{amsmath}
\usepackage{amssymb}
\usepackage{stellar}
\usepackage{definitions}

\begin{document}

\id{psicologia-salute-benessere}
\genpage

\section{Salute e benessere}

\begin{snippetdefinition}{salute-definition}{Salute}
    La \emph{salute} è una risorsa per la vita quotidiana e non lo scopo
    dell'esistenza. È un concetto positivo che valorizza le risorse sociali e
    personali, oltre alle capacità fisiche.
\end{snippetdefinition}

\begin{snippet}{salute-oms}
    La costituzione dell'OMS del 1948 definisce la salute come uno stato di
    completo benessere fisico, sociale e mentale, non soltanto come assenza di
    malattia o infermità. Una persona non affetta da malattia non è quindi
    necessariamente in salute.
\end{snippet}

\begin{snippet}{benessere-contesto}
    Benessere fisico e benessere psichico sono condizioni ottimali che devono
    essere considerate in relazione al contesto socioculturale in cui la persona
    vive e si esprime.
\end{snippet}

\begin{snippetdefinition}{psicologia-positiva-definition}{Psicologia positiva}
    La \emph{psicologia positiva} è un approccio allo studio della personalità e
    del benessere che si concentra sulle caratteristiche di coloro che rispondono
    costruttivamente alle difficoltà della vita e mostrano una forte motivazione
    all'autorealizzazione.
\end{snippetdefinition}

\begin{snippetdefinition}{edonismo-definition}{Edonismo}
    L'\emph{edonismo} è una concezione della felicità come stato di appagamento
    connesso al raggiungimento di un equilibrio omeostatico ottenuto attraverso
    il soddisfacimento dei desideri.
\end{snippetdefinition}

\includesnpt[width=62\%,src=/snippet/static-psychology/hedonic-well-being.jpg]{centered-img}

\begin{snippetdefinition}{eudaimonia-definition}{Eudaimonia}
    L'\emph{eudaimonia} è una concezione della felicità come processo di
    costruzione continua, crescita di complessità, autorealizzazione e
    armonizzazione del benessere individuale con il benessere collettivo.
\end{snippetdefinition}

\includesnpt[width=62\%,src=/snippet/static-psychology/eudaimonic-well-being.jpg]{centered-img}

\begin{snippet}{edonismo-eudaimonia-confronto}
    Edonismo ed eudaimonia non identificano la felicità con il possesso di
    ricchezze o beni materiali. Si distinguono però perché l'edonismo è legato
    al soddisfacimento dei desideri e all'affettività positiva, mentre
    l'eudaimonia riguarda crescita personale, apertura, impegno e armonizzazione
    del benessere individuale con quello collettivo.
\end{snippet}

\begin{snippetdefinition}{set-point-definition}{Set point}
    Il \emph{set point} è l'abituale stato affettivo di una persona, modificato
    solo transitoriamente da intense emozioni positive o negative.
\end{snippetdefinition}

\begin{snippetdefinition}{felicita-definition}{Felicità}
    La \emph{felicità} è la capacità di provare emozioni positive.
\end{snippetdefinition}

\begin{snippetdefinition}{ottimismo-definition}{Ottimismo}
    L'\emph{ottimismo} è una predisposizione interiore di fiducia verso ciò che
    può accadere nella vita.
\end{snippetdefinition}

\begin{snippet}{personalita-benessere}
    Gli studi sul benessere chiedono spesso alle persone di riferire vissuti di
    felicità e infelicità e di valutare quanto siano soddisfacenti le loro vite.
    Quattro tratti sembrano caratterizzare le persone felici: sentirsi contenti
    di sé, percepire controllo sulla propria vita, essere ottimisti ed essere
    estroversi.
\end{snippet}

\end{document}
